\section{Introduction}
A reasoning model may verify a derivation, compare explanations, or discard an unoriginal research idea. These activities share operations of checking, exploration, and revision, but organize them for different purposes. This question of organization connects to Minsky's account of interacting simpler processes and flexible task recruitment in human neuroscience \citep{minsky1986society,cole2013hubs}. For language models, we ask how observable operations vary across tasks, and what those differences imply for intervention.

Recent work describes reasoning behaviors and temporal structure \citep{gandhi2025habits,chen2025lot,li2025thinkarm,lee2026reasonops}, connects conversational organization to learning \citep{kim2026society}, and modifies reasoning through steering and reward-trained vectors \citep{chen2025seal,sinii2025vectors}. These results motivate three distinct questions: can a profile \emph{describe} task differences, \emph{prescribe} beneficial changes, or \emph{explain} an intervention's outcomes? A behavior associated with stronger reasoning need not be beneficial to increase within an already capable model.

Thought Atlas addresses these questions through a shared representation. A \emph{reasoning signature} is the prevalence and temporal allocation of annotated operations in emitted text, without assuming exhaustive access to internal computation. The atlas spans seven models and eight benchmark conditions, including moral advice and idea generation. Within-model comparisons test task information beyond length; matched instructions test responsiveness to requested activity on the same problems.

The descriptive results are consistent, but conditional outcome prediction provides no statistically supported increment after correction. We therefore directly evaluate activation interventions derived from domain contrasts. Large short-budget gains and losses accompany changes in answer production, while manipulation checks reveal unequal fidelity. A selected math archetype, saved reward-trained biases, and SEAL provide comparisons whose average annotation shifts resemble each other more than their outcomes do. Open-task evaluation further separates answer availability from judged originality.

The contribution is this empirical separation of description, prescription, and explanation. It identifies what the common representation captures, where that evidence stops, and which additional controls are needed before interpreting profiles as useful control targets.
